\documentclass[11pt]{article}
\usepackage[margin=0.86in]{geometry}
\usepackage{amsmath,amssymb,amsthm,mathtools,microtype}
\usepackage[dvipsnames]{xcolor}
\usepackage[colorlinks=true,linkcolor=MidnightBlue,citecolor=MidnightBlue,urlcolor=MidnightBlue]{hyperref}
\usepackage{enumitem}
\setlist{nosep,leftmargin=1.6em}
\newtheorem{theorem}{Theorem}
\newtheorem{lemma}[theorem]{Lemma}
\newtheorem{remark}[theorem]{Remark}
\newcommand{\B}{\mathcal B}
\newcommand{\K}{\mathcal K}
\newcommand{\MD}{\operatorname{MD}}
\newcommand{\N}{\mathbb N}
\title{A diffuse masa core for Koszmider's non-diagonalizable state}
\author{Autonomous proof attempt for arXiv:2002.05230}
\date{11 August 2026}

\begin{document}
\maketitle

\begin{abstract}
Koszmider asked whether his ZFC non-diagonalizable pure state on
$\B(\ell_2)$ can restrict multiplicatively to a nonatomic masa.  We do not
settle the weak-closure question.  We prove, however, that every pure state
arising from the projection face in his construction is multiplicative on a
weak-operator dense copy of $C([0,1])$ inside a nonatomic masa.  More
precisely, the state's multiplicative domain contains a positive
multiplicity-one contraction with essential spectrum $[0,1]$.  We then
identify the exact remaining obstruction: continuous endpoint evaluation
does not determine a character on diffuse $L_\infty$.
\end{abstract}

\section{Question and result}

Write $H=\bigoplus_{m\ge1}H_m$ for the finite-dimensional tensor-block
decomposition in \cite{Koszmider}.  Its projection family has the form
\[
 P_\alpha=\bigoplus_{m\ge1}P_{\alpha,m},\qquad \alpha\in2^\N,
\]
where $P_{\alpha,m}$ is rank one on the tensor coordinate indexed by
$\alpha|m$ and is the identity on all other coordinates.  If
$\alpha|m\ne\beta|m$, then $P_{\alpha,m}$ and $P_{\beta,m}$ act on distinct
tensor factors and commute.  The construction chooses a pure state $\phi$
from the face
\[
 \mathcal F=\{\rho:\rho\text{ is a state on }\B(H),\quad
                    \rho(P_\alpha)=1\text{ for all }\alpha\in2^\N\}.
\]
It is non-diagonalizable, hence its restriction to no atomic masa is
multiplicative.  The source asks whether it can be multiplicative on a
nonatomic masa \cite[p.~2]{Koszmider}.

\begin{theorem}[Diffuse $C^*$-core]\label{thm:main}
Let $\phi$ be any pure state in $\mathcal F$.  There are a positive
contraction $A\in\B(H)$ and a nonatomic masa $\mathcal M\subset\B(H)$ such
that
\[
 \sigma(A)=\sigma_{\rm e}(A)=[0,1],\qquad
 \mathcal M=W^*(A),\qquad
 \phi(f(A))=f(1)\quad(f\in C([0,1])).
\]
In particular, $C^*(A)\cong C([0,1])$ is weak-operator dense in $\mathcal M$
and $\phi$ is multiplicative on $C^*(A)$.
\end{theorem}

Thus the source state admits a diffuse masa on which all continuous moments
already form a character.  The theorem stops strictly short of asserting
that $\phi|_{\mathcal M}$ is pure.

\section{A commuting Boolean family in the multiplicative domain}

For a state $\rho$, let
\[
 \MD(\rho)=\{x:\rho(x^*x)=|\rho(x)|^2=\rho(xx^*)\}.
\]
This is a unital $C^*$-subalgebra, and $\rho(xy)=\rho(x)\rho(y)$ whenever
$x\in\MD(\rho)$ and $y\in\B(H)$.

\begin{lemma}\label{lem:md}
$\K(H)\subset\MD(\phi)$ and $P_\alpha\in\MD(\phi)$ for every $\alpha$.
\end{lemma}
\begin{proof}
A pure state of $\B(H)$ which does not annihilate $\K(H)$ is a vector state:
its irreducible GNS representation restricts nontrivially to the compacts,
whose unique nonzero irreducible representation is the identity
representation.  A vector state is diagonalizable, whereas $\phi$ is not.
Hence $\phi|_{\K(H)}=0$.  For compact $k$ this gives
$\phi(k^*k)=\phi(kk^*)=0$, so $k\in\MD(\phi)$.  Finally,
$\phi(P_\alpha)=\phi(P_\alpha^2)=1$, which puts every $P_\alpha$ in the same
domain.
\end{proof}

Choose distinct sequences $\alpha_1,\alpha_2,\ldots\in2^\N$.  Inductively
choose increasing integers $N_n$ so that
\[
 \alpha_n|m\ne\alpha_j|m\quad(m\ge N_n,\ j<n),
\]
and truncate only finitely many blocks:
\[
 E_n=0\oplus\cdots\oplus0\oplus\bigoplus_{m\ge N_n}P_{\alpha_n,m}.
\]
Then $E_n-P_{\alpha_n}$ has finite rank.  Lemma \ref{lem:md} gives
$E_n\in\MD(\phi)$ and $\phi(E_n)=1$.

\begin{lemma}\label{lem:boolean}
The projections $E_n$ commute.  Moreover, for finite disjoint sets
$F,G\subset\N$, the signed Boolean word
\[
 R_{F,G}=\prod_{n\in F}E_n\prod_{n\in G}(I-E_n)
\]
has infinite rank.
\end{lemma}
\begin{proof}
For $i<j$, the $j$th projection is zero before block $N_j$; after that block
the two source projections occupy distinct tensor coordinates.  Thus they
commute everywhere.

Fix finite disjoint $F,G$.  On every sufficiently late block, all relevant
binary prefixes are distinct.  The compression of $R_{F,G}$ to that block is
the tensor product of rank-one projections for indices in $F$, nonzero
complements of rank-one projections for indices in $G$, and identities on
unused coordinates.  It is nonzero.  This happens on infinitely many
orthogonal blocks, so $R_{F,G}$ has infinite rank.
\end{proof}

Set
\[
 D=\sum_{n=1}^\infty 2^{-n}E_n.
\]
The series converges in norm inside $\MD(\phi)$ and $\phi(D)=1$.  Lemma
\ref{lem:boolean} says that the essential joint spectrum of $(E_n)$ is the
full binary cube: every finite cylinder has a nonzero image in the Calkin
algebra.  Applying the binary expansion map gives
\begin{equation}\label{eq:spectrum}
 \sigma(D)=\sigma_{\rm e}(D)=[0,1].
\end{equation}
(Equivalently, every dyadic cylinder has infinite-dimensional spectral
range, and their mesh tends to zero.)

\section{From an atomic model to a diffuse generator}

Let $T$ be multiplication by $t$ on $L_2([0,1])$.  It is a positive
multiplicity-one contraction and
\[
 \sigma(T)=\sigma_{\rm e}(T)=[0,1],\qquad
 W^*(T)=L_\infty([0,1]),
\]
the standard nonatomic masa.  We use the following classical form of the
Weyl--von Neumann theorem: two bounded self-adjoint operators with the same
essential spectrum and the same multiplicities at isolated spectral points
off that spectrum are unitarily equivalent modulo a compact operator
\cite{Berg}.  Here there are no isolated points off the essential spectrum.
By \eqref{eq:spectrum}, some unitary $U:L_2([0,1])\to H$ and compact $K$ obey
\[
 A:=UTU^*=D+K.
\]
Lemma \ref{lem:md} and the $C^*$-algebra property of the multiplicative domain
give $A\in\MD(\phi)$.  Therefore
\[
 \phi(A)=\phi(D)+\phi(K)=1,
 \qquad \phi(f(A))=f(1)\quad(f\in C([0,1])).
\]
Finally $W^*(A)=U L_\infty([0,1])U^*$ is a nonatomic masa and $C^*(A)$ is
weak-operator dense in it.  This proves Theorem \ref{thm:main}.

\section{Why the full question remains}

Let $e_\varepsilon=1_{[0,1-\varepsilon]}(A)\in W^*(A)$.  Since
$I-A\ge\varepsilon e_\varepsilon$ and $\phi(I-A)=0$, positivity yields
\[
 \phi(e_\varepsilon)=0,
 \qquad
 \phi(1_{(1-\varepsilon,1]}(A))=1
 \quad(\varepsilon>0).
\]
Thus $\phi|_{W^*(A)}$ is concentrated in every neighborhood of the endpoint
$1$.  But this does not force a character on diffuse $L_\infty$: endpoint
evaluation on $C([0,1])$ has many character extensions to $L_\infty$, and a
nontrivial convex combination of two such extensions has the same values on
all continuous functions and on all endpoint neighborhoods while failing to
be multiplicative.

The missing step is therefore not spectral or compact-perturbative.  It is a
complete Boolean-lattice statement about all measurable spectral projections.
Continuous functional calculus survives modulo compacts; internal spectral
projections are discontinuous at points of $[0,1]=\sigma_{\rm e}(A)$ and are
not controlled by Weyl--von Neumann equivalence.  This is the precise gap
between the theorem above and the source's nonatomic-masa question.

\begin{thebibliography}{9}
\bibitem{Koszmider}
P. Koszmider, \emph{A non-diagonalizable pure state}, Proc. Natl. Acad. Sci.
USA \textbf{117} (2020), 33084--33089; arXiv:2002.05230v3.

\bibitem{Berg}
I. D. Berg, \emph{An extension of the Weyl--von Neumann theorem to normal
operators}, Trans. Amer. Math. Soc. \textbf{160} (1971), 365--371.
\end{thebibliography}

\end{document}
